\documentclass[11pt]{article}
\usepackage{amsmath,amssymb,amsthm,hyperref}
\newtheorem{theorem}{Theorem}
\newtheorem{lemma}[theorem]{Lemma}
\title{An explicit threshold for eventual equal-sided fair sizes}
\author{Quantitative companion to the positive-completion theorem}
\date{30 September 2026}
\begin{document}
\noindent\textbf{Later improvement.} The signed-target compression below
remains an input to \path{factorial_order_eventual_conductor.tex}, which
improves the eventual-size threshold to $\exp(O(n\log n))$ using a new
prefix completion. The older conductor estimate in this note is retained
as an independently proved bound.

\maketitle

The qualitative positive-completion and Ore argument is due to the
parallel research note in \texttt{../positive\_saturation/}. This companion
extracts bounds from that argument and from
\texttt{signed\_fair\_signatures.tex}. An independent mathematical
review of both the majorant and the base-digit compression is complete;
this is not machine-checked formal verification.

\begin{theorem}\label{main}
Put $P=\prod_{p\le n}p$, $Q=(n!)^{n-1}$, and, for $n\ge2$, define
\[
 L_*=2^{2n+3}n^4n!P,\qquad q=\lceil\log_2L_*\rceil,
 \qquad B_n=Q2^{nq}.
\]
Every positive integer $M$ divisible by $P$ and satisfying
$M\ge B_n^2/P$ is the common number of faces of an $n$-die
permutation-fair collection. In particular there is a threshold
\[
 M_0(n)\le \exp\bigl(O(n^2\log n)\bigr)
\]
above which subset-factorial divisibility is the only obstruction to
equal-sided sizes.
\end{theorem}

This is a bound on the threshold beyond which \emph{every} admissible
size occurs. It does not improve the smallest-size upper bound
$2^{25}(n!)^2$.

Use the reduced squarefree associative algebra from the signed-signature
note, and write $T_a=\exp(a\sum_iX_i)$. A positive word has ordinary
length equal to its total number of letters.

\begin{lemma}[Quantitative positive completion]\label{completion}
Every positive word $u$ of length at most $H$, where $H\ge1$, is a
prefix of a positive word with signature $T_A$, for some integer
$1\le A\le QH$.
\end{lemma}
\begin{proof}
First append letters until all counts are equal to
$a=\max(1,\max_i c_i(u))\le H$. If the logarithm is
$a\sum_iX_i+D_k+\text{higher degrees}$, concatenate all $n!$
relabelings, placing the identity relabeling first. The degree-$k$
term is $\sum_{\sigma\in S_n}\sigma D_k$, since the common degree-one
parts commute and their brackets with $D_k$ have degree at least $k+1$.
This sum vanishes: an invariant multilinear Lie polynomial of degree
$k>1$ is zero. Indeed, on each $k$-letter support invariance forces all
associative-word coefficients to be equal, whereas abelianization of
a Lie polynomial of degree greater than one vanishes. Concatenation
thus kills the degree-$k$ error and multiplies all counts by $n!$.
Apply this for $k=2,\ldots,n$, including a harmless averaging step
when the current error vanishes. The original word remains a prefix,
and the final common count is exactly $A=Qa$.
\end{proof}

\begin{lemma}[Short common right multiples: Mal'cev's identity]\label{growth}
Two positive words of length at most $H\ge1$ have a common right
multiple of length at most $2^nH$.
Each of the two right complements can be chosen with length at most
$(2^n-1)H$.
\end{lemma}
\begin{proof}
The group of reduced signatures is nilpotent of class at most $n$:
if $\mathfrak m$ is the positive-degree ideal, then
$[1+\mathfrak m^a,1+\mathfrak m^b]\subseteq1+\mathfrak m^{a+b}$ and
$\mathfrak m^{n+1}=0$.
Given words $a,b$, define positive words
\[
 A_0=a,\quad B_0=b,\qquad
 A_{j+1}=A_jB_j,\quad B_{j+1}=B_jA_j.
\]
Writing $g_j=A_jB_j^{-1}$ in the group, one has
\[
 g_{j+1}=[A_j,B_j]=[g_j,B_j].
\]
Inductively $g_j\in1+\mathfrak m^{j+1}$, so $A_n=B_n$ as
signatures. Both contain exactly $2^n$ original blocks; they begin
with $a$ and $b$, respectively. Removing those first blocks gives
the stated positive complements.

This is the classical positive nilpotency identity of Mal'cev,
recalled in \cite[\S5.1]{grigorchuk}; its use here requires no
general semigroup theorem.
\end{proof}

\begin{lemma}[Short right fractions]\label{fraction}
A signed word of length $L\ge1$ can be written as $uv^{-1}$, with
$u,v$ positive words of length at most $2^{nq}$, where
\[
 q=\lceil\log_2L\rceil.
\]
Consequently its group element $g$ satisfies $gT_A\in\mathcal M_+$
for some positive fair size $A\le Q2^{nq}$, where $\mathcal M_+$ is
the monoid of signatures of positive words.
\end{lemma}
\begin{proof}
Each signed generator has a right-fraction representation whose
numerator and denominator have length at most one. Suppose two
fractions $u_1v_1^{-1}$ and $u_2v_2^{-1}$ have all four lengths at
most $H$. Lemma~\ref{growth} gives $v_1a=u_2b$ with complements
of length at most $(2^n-1)H$. Their product is
\[
 (u_1a)(v_2b)^{-1}.
\]
Both new lengths are at most $2^nH$.
A balanced binary tree, padded by empty fractions if necessary, has
at most $q$ levels and proves the bound.
Finally complete $vc=T_A$ by Lemma~\ref{completion};
then $gT_A=uc$ is positive and $A\le Q2^{nq}$.
\end{proof}

\begin{lemma}[A short signed fair target]\label{signedlength}
The integral target $T_P$ has a signed-word representative of length
at most $L_*$.
\end{lemma}
\begin{proof}
The signed-signature theorem applies because $P$ satisfies every
subset-factorial divisibility condition. We bound its constructive
degree-correction proof. Start with
$g_1=\prod_i(1+X_i)^P$, of length $nP$. Write
$A=2nP$. For the correction in degree $k$, let $b_{S,\pi}$ be the
integer coefficients in the fixed-last-letter multilinear Lie basis,
and put
\[
 c_k=2^{k-1}\sum_{S,\pi}|b_{S,\pi}|.
\]
Each basis bracket has associative coefficient $\ell^1$ norm
$2^{k-1}$. Its coordinate is a coefficient of a distinct associative
word, so the sum of absolute basis coordinates is no larger than the
coefficient $\ell^1$ norm of the degree-$k$ error.

Products of correction factors and their inverses are coefficientwise
majorized, in these norms, by $\exp(\sum c_jx^j)$. Both the initial
inverse $g_1^{-1}$ and the target $T_P$ are majorized by
$\exp(nPx)$. Therefore
\[
 c_k\le2^{k-1}[x^k]
       \exp\left(Ax+\sum_{2\le j<k}c_jx^j\right).
\]
Inductively,
\[
 c_k\le 2^{k^2}A^k.\tag{*}
\]
To check this, normalize $d_j=c_j/A^j$, set $d_1=1$, and expand the
exponential over partitions of $k$ into parts less than $k$. There are
at most $2^{k-1}$ partitions. In each one,
$\sum_j j^2m_j\le(k-1)^2+1=k^2-2k+2$, since it has at least two
positive parts. The factorial denominators $\prod m_j!$ only decrease
the bound. Thus the recurrence is at most
$2^{k-1}2^{k-1}2^{k^2-2k+2}=2^{k^2}$, proving (*).

It remains to represent a potentially large coefficient efficiently.
For a coefficient $b=b_{S,\pi}$ choose the integer
$R=2^kA+1$, so $|b|<R^k$ by (*). Write
$|b|=\sum_{j=0}^{k-1}d_jR^j$ in base $R$.
Each nonzero term $d_jR^j$ is a product of $k$ positive integers at
most $R$: use $d_j$, $j$ copies of $R$, and $k-j-1$ copies of $1$.
Put the sign of $b$ into the first factor.

Replace the $k$ generators in a nested group commutator by those
signed integer powers. The result is exactly
$1+\operatorname{sgn}(b)d_jR^j B_{S,\pi}$, since repeated-letter
products vanish. Its signed-word length is at most
$(3\cdot2^{k-1}-2)R\le k2^kR$.
Multiplying the at most $k$ digit factors adds their coefficients,
because $B_{S,\pi}^2=0$. Thus $1+bB_{S,\pi}$ costs at most
$k^2 2^kR\le k^2 2^{2k+1}A$ letters.

There are $\binom nk(k-1)!\le n!$ coordinates in degree $k$.
The complete signed word therefore has length
\[
 L\le nP+\sum_{k=2}^n n! k^2 2^{2k+1}A
   \le 2^{2n+3}n^4n!P=L_*.
\]
\end{proof}

\begin{proof}[Proof of Theorem~\ref{main}]
Apply Lemmas~\ref{fraction} and~\ref{signedlength} to $g=T_P$.
There is a positive fair word with common count $A\le B_n$, and
$T_PT_A=T_{A+P}$ also has a positive representative. Arithmetic
necessity implies $P\mid A$. Set $a=A/P$. Positive concatenation
therefore realizes every nonnegative integer combination of $a$ and
$a+1$, multiplied by $P$. Every integer $t\ge a(a-1)$ is such a
combination: write $t=qa+r$, $0\le r<a$; then $q\ge a-1\ge r$ and
$t=(q-r)a+r(a+1)$. Consequently every admissible size at least
$A(A-P)/P\le B_n^2/P$ is realized.

Finally even the elementary bound $P\le n!$ gives
$\log L_*=O(n\log n)$, while $\log Q=O(n^2\log n)$.
The definition of $B_n$ yields $\log B_n=O(n^2\log n)$, as asserted.
\end{proof}
\begin{thebibliography}{9}
\bibitem{grigorchuk} R.~I.~Grigorchuk,
\emph{On Growth in Group Theory}, Proceedings of the International
Congress of Mathematicians, Kyoto 1990, Vol.~I, Mathematical Society
of Japan, 1991, pp.~325--338, especially \S5.1, printed p.~331.
\href{https://www.mathunion.org/fileadmin/ICM/Proceedings/ICM1990.1/ICM1990.1.ocr.pdf}{Primary proceedings scan}.
The positive identity is attributed there to A.~I.~Mal'cev,
\emph{Nilpotent semigroups}, Uchen. Zap. Ivanovsk. Ped. In-ta
4 (1953), 107--111.
\end{thebibliography}
\end{document}
